\documentclass{article}
\usepackage[T1]{fontenc}
\usepackage{lmodern}
\usepackage{graphicx}
\usepackage{amsmath,amssymb}
\usepackage[a4paper,margin=2cm]{geometry}
\usepackage[absolute,overlay]{textpos}
\setlength{\TPHorizModule}{1mm}
\setlength{\TPVertModule}{1mm}
\usepackage[indonesian]{babel}
\usepackage{hyperref}
\setlength{\emergencystretch}{2em}
% unit: Halaman 113

\begin{document}

% === Halaman 113 (python/native) ===

113

• Model mesin enigma ada yang menggunakan 3 rotor atau 4 rotor, setiap rotor 

melakukan operasi substitusi cipher abjad-tunggal. 

• Untuk mesin enigma 4-rotor, berarti terdapat 26  26  26  26 = 456.976 

kemungkinan huruf cipherteks sebagai pengganti huruf plainteks sebelum terjadi 

perulangan urutan cipherteks. 

• Setiap kali sebuah huruf selesai disubstitusi, rotor pertama bergeser satu huruf. 

• Setiap kali rotor pertama selesai bergeser 26 kali, rotor kedua bergeser satu huruf. 

Setelah rotor kedua bergeser 26 kali, rotor ketiga bergeser satu huruf. Setelah 

rotor ketiga bergeser 26 kali, rotor keempat bergeser satu huruf.

\end{document}
